\section{Funzionalità principali}
\subsection{Utente non registrato}
\paragraph{Visualizzazione di animali}
Un utente non registrato può navigare liberamente tra le pagine del sito web e visualizzare gli animali disponibili per l'adozione. Può utilizzare i filtri di ricerca per trovare animali in base a criteri specifici come tipo (Cane o Gatto), razza, età, taglia, sesso ma anche filtrare per nome.
\paragraph{Gestione preferiti}
Per ogni utente non amministratore, registrato o meno, è possibile inserire tra i preferiti gli animali di interesse, in modo da poterli ritrovare facilmente in futuro.\\
La particolarità è che i preferiti sono salvati tramite cookie nel browser dell'utente e al momento dell'accesso vengono caricati automaticamente nel database. Nella sezione preferiti, senza accesso, si vede un chiaro messaggio che invita alla registrazione per salvare i preferiti in modo permanente e quindi sincronizzarli con tutti i dispositivi.
\paragraph{Visualizzazione di eventi}
Nella pagina dedicata agli eventi, un utente può visualizzare gli eventi organizzati dal rifugio, con dettagli come data, ora, luogo e descrizione dell'evento.
Alcune delle funzionalità hanno scopi precisi per rispondere a bisogni di specifiche tipologie di utente. Più dettagli nella sezione relativa all'attrazione degli utenti \ref{sec:attrazione_utenti}.
\subsection{Utente registrato}
Oltre a tutte le funzionalità disponibili per gli utenti non registrati, gli utenti registrati hanno specifiche funzionalità aggiuntive.
\paragraph{Effettuare richieste di adozione}
Compilando il form di adozione, presente comodamente nella pagina di dettaglio di ogni animale disponibile, un utente registrato può inviare una richiesta di adozione per l'animale di interesse.
\paragraph{Gestione profilo}
Gli utenti registrati hanno a disposizione una pagina profilo in cui possono visualizzare e modificare le proprie informazioni personali, come nome, cognome, password ma anche immagine profilo.\\
Più interessante è la possibilità di gestire l'account, cambiando email e password, e la possibilità di eliminare l'account in modo permanente.
\paragraph{Gestione richieste di adozione}
Nell'area personale un utente può visualizzare lo stato delle proprie richieste, vedere gli ultimi avvisi di movimenti effettuati sugli animali interessati (in adozione o in valutazione) o quelle effettuate dall'amministratore che lo hanno coinvolto.\\ 
In ultimo l'utente può visionare una richiesta inviata in dettaglio, con tutte le informazioni inserite e annullare la richiesta di adozione se ancora in stato di "in valutazione" o "nuova".
\subsection{Amministratore}
Un utente amministratore può visualizzare il sito web come un utente normale, ma ha accesso a funzionalità aggiuntive accessibili nell'area amministratore.\\
Ci sono tuttevia delle differenze apportate al sito web normale, sia per migliorare le funzionalità amministrative, sia per limitare l'accesso a funzionalità non necessarie.
Per esempio, la pagina Animali mostra un pulsante che reindirizza alla pagina di amministratore "Assegnati a te" che fa riferimento a tutti gli animali attivi (non adottati) che sono stati assegnati all'amministratore. Inoltre è stato rimossa la pagina dei preferiti e la funzionalità di aggiunta ai preferiti poichè pensata esclusivamente per gli utenti senza privilegi di amministratore.
Infine, come è logico che sia, l'amministratore non può inviare richieste di adozione e per questo non ha un'area dedicata.
\paragraph{Area personale e gestione profilo}
Ogni amministratore ha una propria area personale che mostra le attività da completare e le proprie statistiche personali.\\
Inoltre, come per gli utenti normali, l'amministratore ha una pagina di profilo in cui può visualizzare e modificare le proprie informazioni personali. Tuttavia, rispetto un utente normale, all'amministratore non è concessa la modifica (e l'eliminazione) delle informazioni di accesso (email e password). Si presume che le credenziali di accesso siano gestite in modo più sicuro e controllato.